%% Henkin completeness, compactness, and cardinality for Part 2.
%% Items P2.2.1--P2.2.12 form the Tier-1 Henkin method derivation.
%% Later migrated Tier-2 proofs are compressed individually to 40--60%
%% of their v1.1 detex/[A-Za-z]+ source lengths.

\begin{annexitem}{annex:p2.2.1}{thm:p02-first-order-mcs-closure}{Closure of first-order maximal consistent sets}
\begin{proof}
If \(\Delta\derives\varphi\) while \(\varphi\notin\Delta\), maximality
makes \(\Delta\cup\{\varphi\}\) inconsistent.  Consistency forces an open
\(\varphi\)-occurrence; discharging them all yields
\(\Delta\derives\neg\varphi\), contradicting consistency.
Thus \(\Delta\) is deductively closed.  Consistency excludes both formulas;
if \(\varphi\) is absent, maximality and discharge put \(\neg\varphi\) in
\(\Delta\).
Conjunction uses introduction and elimination.  A disjunction with neither
disjunct contradicts their negations.  For implication, modus ponens proves
necessity; either \(\neg\varphi\) or \(\psi\) propositionally yields
\(\varphi\to\psi\), and closure proves sufficiency.
\end{proof}
\end{annexitem}

\begin{annexitem}{annex:p2.2.2}{thm:p02-quantifier-duality}{Classical quantifier duality}
\begin{proof}
From \(\neg\exists x\,\varphi\), assume \(\varphi[y/x]\), where \(y\) is
absent from \(\varphi\) and every open assumption.  Existential introduction
contradicts it; negation introduction, universal
introduction, and alpha-conversion give \(\forall x\,\neg\varphi\).  The
reverse implication eliminates an existential against a universal instance.

Likewise, each witness to \(\exists x\,\neg\varphi\) contradicts
\(\forall x\,\varphi\).  Finally, \(\neg\exists x\,\neg\varphi\) gives
\(\forall x\,\neg\neg\varphi\) by the first duality; double-negation
elimination at a similarly fresh \(y\), followed by universal introduction
and alpha-conversion, contradicts \(\neg\forall x\,\varphi\).  Reductio
completes the last direction; all eigenvariable conditions hold.
\end{proof}
\end{annexitem}

\begin{annexitem}{annex:p2.2.3}{thm:p02-fresh-constant}{Fresh-constant lemma}
\begin{proof}
Only finitely many assumptions \(\Gamma_0\subseteq\Gamma\) occur.  Choose
\(y\) absent from \(\Gamma_0,\varphi\), and the entire derivation, and replace
\(c\) by \(y\).  Induction preserves every rule: \(c\) binds nothing and
occurs in no open assumption, while freshness preserves eigenvariable
conditions.  Thus \(\Gamma_0\derives\varphi[y/x]\).  Universal introduction
and alpha-conversion yield \(\forall x\,\varphi\); the same derivation works
from \(\Gamma\).
\end{proof}
\end{annexitem}

\begin{annexitem}{annex:p2.2.4}{thm:p02-witness-extension}{Witness-extension lemma}
\begin{proof}
Consistency forces nonvacuous use of \(H_{x,\varphi,c}\); \(\neg I\) gives
\(\Sigma\derives\neg H_{x,\varphi,c}\), hence
\(\exists x\,\varphi\) and \(\neg\varphi[c/x]\).  Fresh-constant plus duality
yields \(\forall x\,\neg\varphi\), \(\neg\exists x\,\varphi\),
contradicting consistency.
\end{proof}
\end{annexitem}

\begin{annexitem}{annex:p2.2.5}{thm:p02-new-constant-elimination}{Eliminating newly added constants}
\begin{proof}
Only finitely many new constants \(d_1,\ldots,d_n\) occur in the given
derivation.  Choose distinct variables \(z_1,\ldots,z_n\) absent from every
formula and rule annotation, and replace the \(d_i\) simultaneously.
Induction preserves every inference: replacement commutes with term formation,
connectives, and capture-avoiding substitution; freshness prevents capture
and preserves the original eigenvariable conditions, alpha-equivalence, and
equality inferences.  Since assumptions and conclusion lie in \(\Lang\), they
contain no \(d_i\) and remain unchanged.  The transformed tree is the
required \(\Lang\)-derivation.
\end{proof}
\end{annexitem}

\begin{annexitem}{annex:p2.2.6}{thm:p02-transfinite-lindenbaum}{Transfinite Lindenbaum lemma}
\begin{proof}
Set \(\Gamma_0=\Sigma\); at a successor, adjoin
\(\varphi_\alpha\) exactly when consistency is preserved; at a limit take
the union.  Transfinite induction shows the sequence increasing and every
stage consistent.  At a limit \(\beta\), a
derivation of \(\bottom\) uses a finite set \(F\) of assumptions.  Each
member enters before \(\beta\), and the maximum of those finitely many
stages (stage \(0\) if \(F\) is empty) already contains \(F\), contradicting
that earlier stage's consistency.

Put \(\Delta=\Gamma_\lambda\).  It is consistent and contains \(\Sigma\).
If \(\varphi\notin\Delta\), choose \(\alpha<\lambda\) with
\(\varphi_\alpha=\varphi\).  The successor clause could have rejected
\(\varphi\) only because \(\Gamma_\alpha\cup\{\varphi\}\) was
inconsistent.  The same finite derivation works from the superset
\(\Delta\cup\{\varphi\}\), proving maximality.  Once the enumeration is
fixed, every recursive choice is determined.
\end{proof}
\end{annexitem}

\begin{annexitem}{annex:p2.2.7}{thm:p02-henkin-lindenbaum}{Henkin--Lindenbaum lemma}
\begin{proof}
Let \(\kappa=\max(\abs{\Lang},\aleph_0)\), and add constants
\(\langle d_\xi:\xi<\kappa\rangle\).  The original \(\Sigma\) remains
consistent by new-constant elimination.  Under Choice, syntax counting and
the distinct formulas \(\exists x_0(x_0=d_\xi)\) give bijective
enumerations, of length \(\kappa\), of the existential formulas and later of
all formulas of the expansion.

Recursively, for the \(\alpha\)-th existential
\(E_\alpha=\exists x_{n_\alpha}\varphi_\alpha\), add
\(H_{x_{n_\alpha},\varphi_\alpha,c_\alpha}\), choosing the least indexed
\(d_\xi\) absent from the current theory and \(E_\alpha\).  Before
\(\alpha\), fewer than \(\kappa\) finite witness formulas have appeared,
and \(E_\alpha\) is finite, so a constant remains.  Witness extension
preserves consistency at successors; at limits, a finite inconsistent
derivation would occur earlier.  Call the union \(\Sigma^H\).

Apply the transfinite Lindenbaum lemma to the fixed enumeration of all
expanded formulas, obtaining a maximal consistent
\(\Delta\supseteq\Sigma^H\).  If \(\exists x\,\varphi\in\Delta\), its
designated witness implication is also in \(\Delta\); implication
elimination and deductive closure put the designated instance in \(\Delta\).
Thus \(\Delta\) has the witness property.  Choice was used only for the two
cardinal enumerations and their bookkeeping.
\end{proof}
\end{annexitem}

\begin{annexitem}{annex:p2.2.8}{thm:p02-equality-congruence}{Equality congruence}
\begin{proof}
Equality introduction gives reflexivity; derived symmetry plus equality
elimination on \(r=x\) gives transitivity.  Repeated equality elimination proves
function and relation compatibility.  Closure puts replacements in
\(\Delta\); symmetry gives the reverse relation implication.
\end{proof}
\end{annexitem}

\begin{annexitem}{annex:p2.2.9}{thm:p02-term-model-evaluation}{Term evaluation in the term structure}
\begin{proof}
Structural induction: variables use \(s_0\); functions use
interpretations and the induction hypotheses, including nullaries.
\end{proof}
\end{annexitem}

\begin{annexitem}{annex:p2.2.10}{thm:p02-quantifier-membership}{Quantifier membership}
\begin{proof}
For existentials, witnessing proves forward; introduction and closure prove
reverse.  Universal elimination proves forward.  If a universal is absent,
maximality and duality give \(\exists x\,\neg\varphi\); its witness
\(\neg\varphi[c/x]\) excludes \(\varphi[c/x]\) by consistency, so not every
instance belongs.
\end{proof}
\end{annexitem}

\begin{annexitem}{annex:p2.2.11}{thm:p02-first-order-truth-lemma}{First-order truth lemma}
\begin{proof}
Use strong induction on complexity.  Term evaluation turns equality truth
into equality of classes, equivalently derivability and, by closure,
membership.  Relations use their interpretation, \(\bottom\) consistency,
and Booleans the maximal-consistency clauses.

For every term \(t\), the substitution lemma gives
\[
 \mathcal M_\Delta,s_0[x\mapsto[t]]\entails\varphi
 \quad\Leftrightarrow\quad
 \mathcal M_\Delta,s_0\entails\varphi[t/x].
\]
The instance has the matrix's complexity, strictly below its quantified
formula.  Induction and quantifier membership handle ``some'' and ``every'';
because every domain element is \([t]\), both semantic quantifiers are
exhausted.
\end{proof}
\end{annexitem}

\begin{annexitem}{annex:p2.2.12}{thm:p02-fol-completeness}{Strong completeness of classical first-order logic}
\begin{proof}
Argue contrapositively.  If \(\Gamma\notderives\varphi\), then
\(\Sigma=\Gamma\cup\{\neg\varphi\}\) is consistent.  Otherwise a finite
derivation of \(\bottom\) either uses no open \(\neg\varphi\), when
explosion derives \(\varphi\) from \(\Gamma\), or discharges all its open
occurrences by reductio to the same conclusion.

Henkin--Lindenbaum extends \(\Sigma\) to a maximal consistent witness set
\(\Delta\).  The truth lemma makes its term structure under \(s_0\) satisfy
every member of \(\Gamma\) and \(\neg\varphi\).  Reduct invariance preserves
these facts in the reduct to \(\Lang\), which is therefore a
\(\Lang\)-countermodel to \(\Gamma\entails\varphi\).
\end{proof}
\end{annexitem}

\begin{annexitem}{annex:p2.2.13}{thm:p02-fol-completeness-countable}{Completeness for countable signatures without choice}
\begin{proof}
The arbitrary-signature proof uses Choice only for cardinal bookkeeping and
the two enumerations in Henkin--Lindenbaum; transfinite Lindenbaum itself is
choice-free once its enumeration is supplied.  Here both inputs are explicit.
Fix one injection of the signature symbols into \(\N\), which selects one
witness rather than a family, and add constants \(d_0,d_1,\ldots\).
Interleave the finite logical alphabet and ordered variables, the signature
symbols under that injection, and the \(d_k\).  Iterated finite-string pairing
assigns each expanded formula a distinct natural-number code.  Least-number
recursion through the code set gives a bijection with all formulas; restricting
and again taking least unused codes enumerates the existential formulas.

Run the witness and Lindenbaum recursions along \(\omega\) and
\(\omega+1\).  Before stage \(n\), only finitely many witness implications,
hence finitely many \(d_k\), have appeared, so the least fresh one exists;
the cardinal-arithmetic appeals reduce to this finite observation and the
displayed coding.  No later choice is hidden: new-constant elimination and
witness extension transform single finite derivations; equality congruence
proves representative independence; \(s_0(x)=[x]\) is explicit; and the
truth and reduct lemmas are inductions.  The same contrapositive construction
therefore proves completeness without Choice.
\end{proof}
\end{annexitem}

\begin{annexitem}{annex:p2.2.14}{thm:p02-fol-lk-completeness}{Completeness of the first-order sequent calculus}
\sketchnote{the individual natural-deduction-to-sequent translation trees and
the alpha-conversion and eigenvariable subinductions are omitted; the route
uses cut, and no cut-elimination or cut-free completeness proof is claimed.}
Soundness is the registered first-order sequent soundness theorem.  For the
converse, induction on a natural-deduction derivation translates propositional
steps as in the classical propositional proof; translates quantifier and
equality steps with their registered sequent counterparts; and uses cut for
eliminations.  Separate subinductions preserve bound-variable
alpha-equivalence and every eigenvariable freshness condition.  Thus a finite
\(\Gamma_0\derives\varphi\) yields
\(\Gamma_0\Rightarrow\varphi\).

Assume the displayed semantic condition and list the support of the succedent
as \(\delta_1,\ldots,\delta_m\), allowing \(m=0\).  The set consisting of
the antecedent support and all \(\neg\delta_j\) is unsatisfiable.  Strong
completeness gives a derivation of \(\bottom\) from it, using finitely many
antecedent formulas and negated succedent formulas; translate that derivation.
Cut with the registered bottom-left initial sequent to empty the succedent.
For each used \(\neg\delta_j\), identity followed by the registered
negation-right rule gives a sequent with \(\delta_j,\neg\delta_j\) on the
right; a cut moves \(\delta_j\) to the current succedent.  Weakening adds
the requested occurrences and contraction removes excess duplicates.
If \(m=0\), no such cuts occur.  Empty antecedents are handled identically.
\end{annexitem}

\begin{annexitem}{annex:p2.2.15}{thm:p02-fol-compactness}{First-order compactness}
\begin{proof}
A model of \(T\) models every finite subset.  Conversely, if all finite
subsets are satisfiable but \(T\) is inconsistent, a derivation of
\(\bottom\) uses finite \(T_0\subseteq T\), contradicting soundness.
The Henkin construction and reduct invariance therefore give a model of
\(T\).  If \(T\entails\sigma\), completeness gives a finite derivation of
\(\sigma\); the finite assumptions \(T_0\) it uses satisfy
\(T_0\entails\sigma\) by soundness.
\end{proof}
\end{annexitem}

\begin{annexitem}{annex:p2.2.16}{thm:p02-finite-structures-not-axiomatizable}{Finite structures are not first-order axiomatizable}
\begin{proof}
If \(T\) axiomatized finite nonempty structures, every finite part of
\(T\cup\{\iota_n:n\geq2\}\) would have a large finite model.  Compactness
would yield a model infinite by the \(\iota_n\) and finite by \(T\).
\end{proof}
\end{annexitem}

\begin{annexitem}{annex:p2.2.17}{thm:p02-tarski-vaught}{Tarski--Vaught test}
\begin{proof}
If \(\mathcal N\preccurlyeq\mathcal M\), an existential true in
\(\mathcal M\) at parameters from \(\mathcal N\) is true in
\(\mathcal N\), hence has a witness there.

Conversely, induct on formulas.  Substructure closure makes term values agree,
so atomic and Boolean cases follow.  An existential true in \(\mathcal N\)
has a witness also in \(\mathcal M\); one true in \(\mathcal M\) has, by
hypothesis, a witness in \(\mathcal N\); induction transfers the matrix both
ways.  For a universal, truth in \(\mathcal M\) transfers every instance at
an element of \(\mathcal N\).  If false in \(\mathcal M\), apply the witness
condition to \(\exists y\,\neg\varphi\) and transfer that counterinstance.
Thus all formulas agree.
\end{proof}
\end{annexitem}

\begin{annexitem}{annex:p2.2.18}{thm:p02-downward-ls}{Downward L\"owenheim--Skolem}
\begin{proof}
Put \(\mu=\max(\abs{A},\abs{\Lang},\aleph_0)\), choose
\(m_0\in\dom\mathcal M\), and use Choice to fix witnesses for every
existential formula and parameter tuple, using \(m_0\) if none exists.  Let
\(A_0\) contain \(A,m_0\), and interpreted constants; obtain \(A_{n+1}\)
by closing \(A_n\) under signature functions and those witness functions.

Syntax counting and finite-tuple coding give \(\abs{A_n}\leq\mu\) by
induction, hence the union \(N=\bigcup_nA_n\) has size at most \(\mu\).
It induces a substructure \(\mathcal N\) containing \(A\).  Parameters of
any true existential lie together in some \(A_n\), whose chosen witness lies
in \(A_{n+1}\); Tarski--Vaught gives
\(\mathcal N\preccurlyeq\mathcal M\).
\end{proof}
\end{annexitem}

\begin{annexitem}{annex:p2.2.19}{thm:p02-upward-ls}{Upward L\"owenheim--Skolem}
\begin{proof}
Add constants \(c_\alpha\) for \(\alpha<\kappa\) and all inequalities
\(c_\alpha\ne c_\beta\).  Each finite fragment expands a fixed infinite
model of \(T\), so compactness gives a model whose named set \(A\) has size
\(\kappa\).  Downward L\"owenheim--Skolem in the expansion gives an
elementary substructure containing \(A\) of size at most \(\kappa\), hence
exactly \(\kappa\); its \(\Lang\)-reduct models \(T\).  Taking
\(T=\operatorname{Th}_{\Lang}(\mathcal M)\) proves the final claim.
\end{proof}
\end{annexitem}

\begin{annexitem}{annex:p2.2.20}{thm:p02-upward-ls-elementary-extension}{Upward L\"owenheim--Skolem via elementary diagrams}
\begin{proof}
Expand \(\Lang(\mathcal M)\) further by constants
\(d_\alpha\) for \(\alpha<\kappa\), and adjoin to
\(\operatorname{ElDiag}(\mathcal M)\) all
\(d_\alpha\ne d_\beta\) for \(\alpha<\beta<\kappa\) and all
\(d_\alpha\ne c_a\) for \(a\in\dom\mathcal M\).  Every finite fragment is
satisfied by the parameter expansion of \(\mathcal M\): because its domain
is infinite, interpret the finitely many occurring \(d_\alpha\) as distinct
elements avoiding the finitely many named \(a\)'s.  Compactness supplies a
model \(\mathcal P\) of the whole theory.

Let \(A\) comprise the interpretations of every \(c_a\) and
\(d_\alpha\).  It has cardinality \(\kappa\).  Downward
L\"owenheim--Skolem in the full expanded signature, whose cardinality is
\(\kappa\), gives
\(\mathcal Q\preccurlyeq\mathcal P\) containing \(A\) and of size at most
\(\kappa\), hence exactly \(\kappa\).  Let \(\mathcal N\) be its
\(\Lang\)-reduct and define \(j(a)=c_a^{\mathcal Q}\).  Distinctness
sentences make \(j\) injective.  For each parameter sentence, it or its
negation belongs to the elementary diagram according to its truth in
\(\mathcal M\), so \(j\) preserves and reflects every \(\Lang\)-formula.
Thus \(j\) is elementary.  The interpretation of
\(d_0\) is outside its range, making the range proper.
\end{proof}
\end{annexitem}

\begin{annexitem}{annex:p2.2.21}{thm:p02-nonstandard-arithmetic-model}{A countable nonstandard elementary model of arithmetic}
\begin{proof}
Apply the preceding theorem to the countable infinite structure
\(\mathcal N_{\mathrm{ar}}\) with \(\kappa=\aleph_0\).  Replacing the target
by an isomorphic copy over the elementary image makes that image an actual
proper substructure.  The resulting countable elementary extension contains
an element outside every embedded standard numeral.
\end{proof}
\end{annexitem}
